%% hopfion-framework/gravity/newton.tex
%% Sources: main_paper7.tex
%% G_N from gradient flux imbalance, Einstein equations, Cassini bound
%% Auto-generated from project files. Edit source papers to update.


%════════════════════════════════════════════════════════════════════
% Newton's Constant and Einstein Equations — Paper VII
%════════════════════════════════════════════════════════════════════
The parent action governing the condensate scalar field $\rho$, gravity
and cosmology is
\begin{equation}
  S[\rho] = \int d^4x\sqrt{-g}\,
  \biggl[\frac{g^{\mu\nu}\partial_\mu\rho\,\partial_\nu\rho}
              {\vp^6(1+\beta\rho)}
         + \Lambda_0\,e^{-\vp^6\rho}\biggr]
  + \int S_{\mathrm{eff}}(\theta,\rho)\,d\Omega,
  \label{P7:eq:parent_action}
\end{equation}
where $S_{\mathrm{eff}}(\theta,\rho) = \sin^4\!\theta/[\vp^6(1+\bst\rho)]$
is the suppression formula of~\cite{Paper1}, a $k$-essence type
scalar-tensor theory in disguise.

\begin{remark}[Extension of the density variable]
\label{P7:rem:rho_extension}
In Papers~I--IV the feedback denominator $1+\bst\kappa$ uses the local
field curvature $\kappa$ (the FN gradient density) as the density variable.
In this paper, the parent action~\eqref{P7:eq:parent_action} uses a scalar
field $\rho$ as an independent dynamical variable, with $1+\beta\rho$ in
the denominator.  The two descriptions are compatible: in the static
condensate background, $\rho\propto\kappa$ and $\beta\propto\bst$, so
Axiom~A2 recovers the Papers~I--IV form.  The $\rho$-based description
is a $k$-essence extension that enables the cosmological and gravitational
analysis of this paper.  All recalled results R1--R5 hold under both
descriptions.
\end{remark}
\status{published}
\source{P7:rem:rho_extension}

\begin{remark}[RG fixed point and the fermion mass tower]
\label{P7:rem:rg_tower}
The $k$-essence scalar-tensor functional~\cite{Paper1} governing the
condensate scalar field $\rho$ is
\begin{equation}
  S_{\mathrm{parent}}[\rho]
  \;=\; \int\!\mathrm{d}^4x\,\sqrt{-g}\;
  \frac{|\nabla\rho|^2}{\vp^6\bigl(1 + \bst\,|\nabla\rho|^2\bigr)},
  \label{P7:eq:parent_kin}
\end{equation}
with $\vp^6 = \lam \approx 17.944$ the Bogomolny parameter.
Under Derrick scaling $r \to \zeta r$, the kinetic sector of
\eqref{P7:eq:parent_kin} scales as $\zeta^{-2}$ and the angular-suppression
sector (from the $\sin^4\!\theta/\vp^6$ factor of the condensate vacuum)
scales as $\zeta^4$.  The unique fixed point of the resulting balance
equation is $\zeta = \vp$, establishing the golden ratio as the
RG scale ratio of the condensate action.

Stable bound states of the condensate field therefore exist at the
discrete scales
\begin{equation}
  m_n \;=\; m_0\,\vp^{-2n}, \qquad n = 0, 1, 2, \ldots,
  \label{P7:eq:mass_spectrum}
\end{equation}
with one level per pair of transverse dimensions in the $\sin^4\!\theta$
suppression.  The full derivation of this tower --- including the proof
that $\zeta = \vp$ is the unique fixed point, the quantisation of levels,
the identification of the ground state with $\Lcond$, and the connection
to the WZW T-matrix phases of Papers~I and~IV --- is carried out in the
companion Paper~VI~\cite{Paper6}.
Equation~\eqref{P7:eq:mass_spectrum} is used throughout this paper as a
derived result; see Paper~VI for all proofs.
\end{remark}
\status{published}
\source{P7:rem:rg_tower}

\begin{remark}[Role in this paper]
\label{P7:rem:role_in_paper}
The mass spectrum~\eqref{P7:eq:mass_spectrum} and the parent
action~\eqref{P7:eq:parent_action} enter the gravitational sector of the
present paper in the following ways:
(i)~The energy-momentum tensor $T_{\mu\nu}[\rho]$ derived from
\eqref{P7:eq:parent_action} sources the modified Friedmann equations of
Section~\ref{P7:sec:de}.
(ii)~The discrete levels~\eqref{P7:eq:mass_spectrum} set the mass scales
of the condensate contributions to the gravitational effective action
of Section~\ref{P7:sec:einstein_eqs}.
(iii)~The tower ground state $m_0 = \Lcond$ fixes the cosmological
constant via the condensate vacuum energy density
(Section~\ref{P7:sec:de}).
The fermion mass predictions themselves (lepton masses, Koide relation,
quark hierarchy) are not re-derived here; they are summarised
in~\cite{Paper1,Paper4,Paper6,Koide1983}.
\end{remark}
\status{published}
\source{P7:rem:role_in_paper}

Integrating out the scalar fluctuation $\xi$ at one loop via the
Seeley--DeWitt heat-kernel expansion gives an effective action whose
$a_2$ coefficient contains the Ricci scalar:
\begin{equation}
  \Gamma_{1\text{-loop}} \supset
  \frac{\LUV^2}{2(4\pi)^2}\int d^4x\sqrt{-g}\,
  \left(\frac{R}{6} - m_\xi^2\right).
  \label{P7:eq:1loop}
\end{equation}
Identifying the induced Planck mass by matching the $R$-coefficient
to the Einstein--Hilbert action $(\MPl^2/2)\int\sqrt{-g}\,R\,d^4x$
confirms that a gravitational sector is induced at the condensate
scale, but does not itself set the Planck mass: $\MPl$ is fixed by the
density-feedback flux, $G_N=G_{\mathrm{src}}^2/(4\pi\vp^6)$
(equation~\eqref{P7:eq:GN}), and~\eqref{P7:eq:1loop} serves as a
consistency check that gravity is induced at the right scale rather
than an independent determination of $\LUV$.
The one-loop action generates a non-minimal coupling
$\xNMC=\LUV^2\beta/[4(4\pi)^2]$~\eqref{P7:eq:xNMC}, and the full
effective action takes the scalar-tensor form
\begin{equation}
  S_{\mathrm{eff}} = \int d^4x\sqrt{-g}\!\left[
    \frac{\MPl^2}{2}R
    + \xNMC\rho R
    + \frac{(\partial\rho)^2}{\vp^6(1+\beta\rho)}
    + \Lambda_{\mathrm{eff}}
  \right].
  \label{P7:eq:Seff}
\end{equation}
The ultraviolet cutoff is set by the condensate's own $\vp$-tower.
Both $\MPl$ and $\LUV$ lie on the tower $M_n=\Lcond\,\vp^{2n}$ of
Paper~VI, separated by exactly three steps:
\begin{equation}
  \boxed{\LUV \;=\; \frac{\MPl}{\vp^6}\;\approx\;0.0557\,\MPl,}
  \label{P7:eq:MPl}
\end{equation}
the gap $\vp^6=\lambda$ being the same Bogomolny/$S_{\mathrm{eff}}$
suppression that fixes $G_N$~\eqref{P7:eq:GN} --- a topological
relation, not a loop coincidence. Numerically
\begin{equation}
  \LUV \;\approx\; 6.8\times10^{17}\,\mathrm{GeV}\;\approx\;\frac{\MPl}{18}.
  \label{P7:eq:LUV_approx}
\end{equation}
Matching~\eqref{P7:eq:Seff} to the Brans--Dicke form gives
\begin{equation}
  \oBD = \frac{\vp^{12}}{3\beta}
  \approx \frac{107.6}{\beta}.
  \label{P7:eq:oBD}
\end{equation}
For $\beta\in[0.1,1.0]$ this gives $\oBD\in[108,1076]$.

\begin{theorem}[Einstein field equations from the k-essence condensate]
\label{P7:thm:einstein}
Varying $S_{\mathrm{eff}}$~\eqref{P7:eq:Seff} with respect to
$g^{\mu\nu}$ yields the exact dynamical Einstein equations
\begin{equation}
  F(\rho)\,G_{\mu\nu}
  \;+\; \xNMC\!\left(\nabla_\mu\nabla_\nu - g_{\mu\nu}\,\square\right)\rho
  \;=\; T_{\mu\nu}^{(\rho)},
  \label{P7:eq:einstein_full}
\end{equation}
where $F(\rho)\equiv\MPl^2/2 + \xNMC\rho$ is the effective gravitational
strength and the soup stress-energy is
\begin{equation}
  T_{\mu\nu}^{(\rho)}
  \;=\; \frac{\partial_\mu\rho\,\partial_\nu\rho}{\vp^6(1+\beta\rho)}
  \;-\; \left[\frac{(\partial\rho)^2}{2\vp^6(1+\beta\rho)}
             + \Lambda_{\mathrm{eff}}\right]g_{\mu\nu}.
  \label{P7:eq:Tmunu}
\end{equation}
In the \emph{frozen attractor limit} $\rho\to\rinfty=\vp/\beta$, $\partial_\mu\rho\to0$:
\begin{equation}
  \boxed{G_{\mu\nu} + \frac{\Lobs}{F(\rinfty)}\,g_{\mu\nu} = 0,}
  \label{P7:eq:einstein_frozen}
\end{equation}
i.e.\ the standard Einstein equation with cosmological constant
$\Lambda_{\mathrm{phys}}=\Lobs/F(\rinfty)$, where the effective
gravitational strength is
\begin{equation}
  F(\rinfty)
  \;=\; \frac{\MPl^2}{2}(1 + 3\vp^7)
  \;=\; \frac{\MPl^2}{2}\,(39\vp + 25).
  \label{P7:eq:Feff}
\end{equation}
The corrections to~\eqref{P7:eq:einstein_frozen} are controlled by
$(\partial_\mu\rho)^2/\Lobs \ll 1$ (kinetic suppression) and
$\xNMC\square\rho/F(\rinfty)\ll1$ (slow-roll suppression),
both satisfied once the field has approached the attractor.
\end{theorem}
\status{published}
\source{P7:thm:einstein}

\begin{remark}[Tree level vs.\ one-loop graviton dynamics]
\label{P7:rem:spin2}
The soup action~\eqref{P7:eq:parent_action} alone does not produce a
propagating spin-2 graviton at tree level: the metric is auxiliary,
determined algebraically by the stress tensor.  Full tensor GR emerges
after one-loop integration, where the
Seeley--DeWitt $a_2$ coefficient induces the Einstein--Hilbert term.
The complete dynamical Einstein equations derived from the resulting
effective action $S_{\mathrm{eff}}$ are proved in Theorem~\ref{P7:thm:einstein}:
the k-essence action reduces exactly to
Einstein--Hilbert plus a cosmological constant in the frozen attractor limit.
\end{remark}
\status{published}
\source{P7:rem:spin2}

The explicit weak-field metric $ds^2=-(1-2GM/r)dt^2+(1+2GM/r)(dr^2+r^2d\Omega^2)$
(the standard isotropic post-Newtonian form) has linearised Ricci
tensor $R_{\mu\nu}=-\tfrac12\square h_{\mu\nu}$ in Lorenz gauge with
$\Phi_N=-GM/r$, giving $\nabla^2\Phi_N=0$ in the exterior, hence
\begin{equation}
  G_{\mu\nu} = 0 \quad\text{in the exterior,}
  \label{P7:eq:Gmunu_zero}
\end{equation}
verified without assuming it.
\status{published}
\source{P7:sec:tensor}

The angular integral $\int S_{\mathrm{eff}}\,d\Omega$ expanded to second
order in $h_{\mu\nu}$ generates a mass term for the transverse-traceless
(TT) modes,
\begin{equation}
  m_g^2 = \frac{16\pi}{315\,\vp^6(1+\beta\rho_\infty)},
  \label{P7:eq:graviton_mass}
\end{equation}
a Lorentz-breaking mass tied to the preferred direction $\hat e_r$
(Axiom~A1) rather than the covariant Fierz--Pauli form. The GW
propagation speed $v_{\mathrm{GW}}^2=1-m_g^2/\omega^2$ gives
$|v_{\mathrm{GW}}-c|/c<10^{-33}$ at LIGO frequencies
($\omega\sim100\,\mathrm{Hz}$, $m_g<10^{-29}\,\mathrm{eV}$), consistent
with the GW170817 bound $|v_{\mathrm{GW}}-c|/c<10^{-15}$. The specific
golden-ratio value $m_g^2\propto\vp^{-8}$ is a Jordan-frame quantity,
tied to the non-canonical $S_{\mathrm{eff}}$ normalisation and
rescaled under the conformal map to the Einstein frame; no observable
depends on it, since only the smallness $m_g\ll\omega$ enters
$v_{\mathrm{GW}}$ and holds for any such $m_g$. At tree level
there is one propagating scalar mode $\xi$ (breathing polarisation) and
no propagating TT modes; including the one-loop induced $R$ term, the
TT modes $h_+,h_\times$ become dynamical, giving three polarisations
in total.
\status{published}
\source{P7:sec:tensor}

\begin{theorem}[Bianchi consistency and energy-momentum conservation]
\label{P7:thm:conservation}
The effective action $S_{\mathrm{eff}}$~\eqref{P7:eq:Seff} is diffeomorphism
invariant, so by Noether's second theorem the contracted Bianchi identity
$\nabla^\mu G_{\mu\nu}=0$ implies a conservation law.
Explicitly:

\noindent\textbf{(i) Scalar field equation.}
Varying $S_{\mathrm{eff}}$ with respect to $\rho$ gives
\begin{equation}
  \nabla_\mu\!\left[\frac{\partial^\mu\rho}{\vp^6(1+\beta\rho)}\right]
  \;+\; \frac{\beta\,(\partial\rho)^2}{2\vp^6(1+\beta\rho)^2}
  \;-\; V'(\rho)
  \;=\; \xNMC\,R,
  \label{P7:eq:scalar_eom}
\end{equation}
where $V'(\rho)=\partial V/\partial\rho=-\vp^6\Lambda_0 e^{-\vp^6\rho}$
and the right-hand side $\xNMC R$ is the NMC source.

\noindent\textbf{(ii) Contracted Bianchi identity.}
On any solution of both~\eqref{P7:eq:einstein_full} and~\eqref{P7:eq:scalar_eom},
\begin{equation}
  \nabla^\mu T_{\mu\nu}^{(\rho)} \;=\; 0,
  \label{P7:eq:Tcons}
\end{equation}
i.e.\ the soup stress-energy is covariantly conserved.

\noindent\textbf{(iii) Equivalence.}
Any two of the three equations
$\{\text{Einstein~\eqref{P7:eq:einstein_full},\;scalar~\eqref{P7:eq:scalar_eom},\;conservation~\eqref{P7:eq:Tcons}}\}$
imply the third.
\end{theorem}
\status{published}
\source{P7:thm:conservation}
\begin{proof}[Proof sketch]
Writing $K=K(\rho)$ and $\mathcal{E}\equiv K\square\rho+K'(\partial\rho)^2/2+V'(\rho)$,
a direct computation gives
\begin{equation}
  \nabla^\mu T_{\mu\nu}^{(\rho)}
  = \mathcal{E}\cdot\partial_\nu\rho,
  \label{P7:eq:div_T}
\end{equation}
using the symmetry of $\nabla_\mu\nabla_\nu\rho$. Taking the divergence
of~\eqref{P7:eq:einstein_full} and using $\nabla^\mu G_{\mu\nu}=0$:
\begin{equation}
  \nabla^\mu\!\bigl[F(\rho)G_{\mu\nu}\bigr]
  = \xNMC(\partial^\mu\rho)G_{\mu\nu}
  = -\,\nabla^\mu\!\bigl[\xNMC(\nabla_\mu\nabla_\nu - g_{\mu\nu}\square)\rho\bigr]
  + \nabla^\mu T_{\mu\nu}^{(\rho)}.
  \label{P7:eq:bianchi_split}
\end{equation}
The NMC divergence evaluates via the Ricci commutation identity
$[\square,\nabla_\nu]\rho=R_\nu{}^\mu\partial_\mu\rho$, giving
$\nabla^\mu T_{\mu\nu}^{(\rho)}=(\mathcal{E}-\xNMC R)\partial_\nu\rho$,
which vanishes on~\eqref{P7:eq:scalar_eom}. (iii) follows since the
derivation of~\eqref{P7:eq:div_T} is reversible.
\end{proof}
\status{published}
\source{P7:thm:conservation}

\begin{remark}[Physical content of Theorem~\ref{P7:thm:conservation}]
\label{P7:rem:conservation_physical_content}
Theorem~\ref{P7:thm:conservation} closes the gravitational sector:
Theorem~\ref{P7:thm:einstein} derives the Einstein equations from $S_{\mathrm{eff}}$,
and Theorem~\ref{P7:thm:conservation} shows that the soup scalar $\rho$ satisfies
an equation of motion whose consistency with gravity is guaranteed by the
diffeomorphism invariance of $S_{\mathrm{eff}}$.
In particular, the stress-energy $T_{\mu\nu}^{(\rho)}$ is automatically
conserved on solutions --- no additional assumption about matter coupling is needed.
The framework is therefore a complete, self-consistent scalar-tensor gravity theory.
\end{remark}
\status{published}
\source{P7:rem:conservation_physical_content}

\begin{theorem}[Cassini resolution via Vainshtein screening~\cite{Vainshtein1972}]
\label{P7:thm:cassini}
All inputs to the Vainshtein radius are fixed by the framework without
free parameters.  Substituting the dark-energy scalar mass
$m_\xi^2 = \vp^{20}\,\Lambda_{\mathrm{obs}}\,\MPl^4$
(equation~\eqref{P7:eq:mxi}, no free parameter) and the Brans--Dicke parameter
$\oBD = \vp^{12}/(3\beta)$ (equation~\eqref{P7:eq:oBD}) into the standard
massive-BD Vainshtein formula~\cite{Vainshtein1972,Dvali2000},
the Vainshtein radius for the Sun is
\begin{equation}
  \boxed{r_V
  \;=\; \left(\frac{r_S}{6\beta\,\vp^8\,
        (\Lambda_{\mathrm{obs}}/\MPl^4)}\right)^{\!1/3}
  \;\approx\; 9.3\times10^6\,\mathrm{AU},}
  \label{P7:eq:rV_exact}
\end{equation}
where $r_S = 2G_N M_\odot/c^2\approx3\,\mathrm{km}$ is the solar
Schwarzschild radius and $\Lambda_{\mathrm{obs}}/\MPl^4\approx10^{-122}$.
This exceeds the solar system by five orders of magnitude.
The screened PPN parameter satisfies
\begin{equation}
  |\gamma_{\mathrm{PPN}}-1|
  \;\approx\;\left(\frac{R_\odot}{r_V}\right)^{1/2}
  \;\approx\; 2.2\times10^{-5}
  \;<\; 2.3\times10^{-5} \quad\text{(Cassini bound).}
  \label{P7:eq:gamma_PPN}
\end{equation}
The Vainshtein nonlinearity condition is satisfied throughout the solar
system: at the Cassini measurement orbit ($r\sim10\,\mathrm{AU}$),
$(r_V/r)^{3/2}\approx10^{9}\gg1$, confirming deep nonlinear
suppression of the fifth force.
\end{theorem}
\status{published}
\source{P7:thm:cassini}
\begin{proof}
$m_\xi^2=\vp^{20}\Lambda_{\mathrm{obs}}\MPl^4$ gives
$m_\xi=\vp^{10}\sqrt{\Lambda_{\mathrm{obs}}}\,\MPl\approx\vp^{10}\sqrt3\,H_0$,
confirming $m_\xi\sim H_0$ with no free parameter. Substituting into
$r_V=(G_NM\,\oBD/m_\xi^2)^{1/3}$ with $\oBD=\vp^{12}/(3\beta)$ and
$G_N=1/\MPl^2$:
\begin{equation}
  r_V^3
  = \frac{M}{\MPl^2}\cdot\frac{\vp^{12}}{3\beta}
    \cdot\frac{1}{\vp^{20}\Lambda_{\mathrm{obs}}\MPl^4}
  = \frac{M}{3\beta\,\vp^8\,\Lambda_{\mathrm{obs}}\,\MPl^6}
  = \frac{r_S}{6\beta\,\vp^8\,(\Lambda_{\mathrm{obs}}/\MPl^4)},
  \label{P7:eq:rV_derivation}
\end{equation}
where $r_S=2M/\MPl^2$ is the Schwarzschild radius. Numerically with
$\beta=0.1$, $r_S=3\,\mathrm{km}$, $\vp^8\approx47$, and
$\Lambda_{\mathrm{obs}}/\MPl^4\approx10^{-122}$:
$r_V\approx9.3\times10^6\,\mathrm{AU}$.
\end{proof}
\status{published}
\source{P7:thm:cassini}

\begin{remark}[Self-consistency of the resolution]
\label{P7:rem:cassini_self}
The resolution is fully self-consistent: $m_\xi\sim H_0$ is \emph{derived}
from the dark energy sector (equation~\eqref{P7:eq:mxi}), not imposed.
The single parameter $\beta$ enters only through $\oBD = \vp^{12}/(3\beta)$;
its framework value $\beta\approx0.1$ gives $r_V\approx9.3\times10^6\,\mathrm{AU}$.
The constraint $|\gamma_{\mathrm{PPN}}-1|<2.3\times10^{-5}$ is satisfied
for $\beta\leq0.13$, consistent with the range of Axiom~A2.
\end{remark}
\status{published}
\source{P7:rem:cassini_self}

\begin{equation}
  G_N = \frac{G_{\mathrm{src}}^2}{4\pi\vp^6}.
  \label{P7:eq:GN}
\end{equation}

\begin{remark}
\label{P7:rem:4pi_factor}
The $4\pi$ factor in~\eqref{P7:eq:GN} comes from the Green's function
normalisation $\nabla^2(1/r)=-4\pi\delta^{(3)}$.
Earlier informal estimates used $G_N=G_{\mathrm{src}}^2/\vp^6$; the
correct formula is~\eqref{P7:eq:GN}.
The UV cutoff from the Seeley--DeWitt one-loop computation (Section~5)
gives $\LUV = \MPl/(4\pi\sqrt{2})\approx0.0563\,\MPl$;
the compact estimate $\LUV\approx\MPl/\vp^6\approx0.0557\,\MPl$
follows from the Bogomolny coincidence $4\pi\sqrt{2}\approx\vp^6$
(see~\eqref{P7:eq:MPl} and~\eqref{P7:eq:LUV_approx}).
\end{remark}
\status{published}
\source{P7:rem:4pi_factor}

\begin{equation}
  G_N^{\rm phys}
  \;=\; \frac{1}{2F(\rinfty)}
  \;=\; \frac{1}{\MPl^2(1+3\vp^7)},
  \label{P7:eq:GN_phys}
\end{equation}

\begin{remark}[Physical interpretation of $F(\rinfty)$ and the renormalised $\MPl$]
\label{P7:rem:MPl_renorm}
The factor $(1+3\vp^7)$ in~\eqref{P7:eq:Feff} arises because the NMC coupling
$\xNMC\rho_\infty$ is not small: at the attractor the soup scalar contributes
$3\vp^7\approx87$ times as much to the effective gravitational strength as the
purely loop-induced $\MPl^2$ term.
The \emph{physical} Newton constant measured at laboratory scales is
therefore the $G_N^{\rm phys}$ above, a factor $(1+3\vp^7)^{-1}\approx1/88$
below the na\"ive one-loop estimate.
This is the self-consistent gravitational coupling of the condensate: both the
loop-induced $\MPl^2/2$ and the attractor NMC term $\xNMC\rinfty$ contribute,
and their precise combination $39\vp+25$ is an exact algebraic consequence of
the golden ratio.
\end{remark}
\status{published}
\source{P7:rem:MPl_renorm}

Varying $\int\!d^4x\sqrt{-g}\,(f(\rho)R)$ with $f(\rho)=\MPl^2/2+\xNMC\rho$
gives
\begin{equation}
  \frac{\delta(\sqrt{-g}\,f R)}{\delta g^{\mu\nu}}
  = \sqrt{-g}\!\left[f R_{\mu\nu} - \tfrac{1}{2}f R g_{\mu\nu}
    + (g_{\mu\nu}\square - \nabla_\mu\nabla_\nu)f\right],
  \label{P7:eq:var_fR}
\end{equation}
so with $\nabla_\mu f=\xNMC\partial_\mu\rho$:
\begin{equation}
  f G_{\mu\nu} + \xNMC(g_{\mu\nu}\square - \nabla_\mu\nabla_\nu)\rho = 0
  \quad[\text{gravity sector}].
  \label{P7:eq:gravity_sector}
\end{equation}
With $\xNMC=\LUV^2\beta/[4(4\pi)^2]$, $\rinfty=\vp/\beta$, and
$\MPl^2/2=\LUV^2/[12(4\pi)^2\vp^6]$:
\begin{equation}
  \frac{\xNMC\,\rinfty}{\MPl^2/2}
  \;=\; \frac{\LUV^2\vp}{4(4\pi)^2}\cdot\frac{12(4\pi)^2\vp^6}{\LUV^2}
  \;=\; 3\vp^7,
  \label{P7:eq:NMC_ratio}
\end{equation}
and therefore
\begin{equation}
  F(\rinfty) \;=\; \frac{\MPl^2}{2}\bigl(1+3\vp^7\bigr),
  \label{P7:eq:Feff_derived}
\end{equation}
which, using the Fibonacci identity $\vp^7=13\vp+8$, gives
$1+3\vp^7=39\vp+25$.

\begin{equation}
  \nabla^2(\delta\rho) = -G_{\mathrm{src}}M\,\delta^{(3)}(\mathbf{r}),
  \label{P7:eq:poisson}
\end{equation}


%════════════════════════════════════════════════════════════════════
% Exact Rotating Condensate Background — Paper VII
%════════════════════════════════════════════════════════════════════
\begin{theorem}[Exact rotating condensate background]
\label{P7:thm:rotating}
Let $\kappa = \beta\vp^2 G_N^{1/2}$.
The unique stationary, axisymmetric, asymptotically-attractor exterior
solution to $\nabla^2\rho_0 = 0$ with the correct source term and
$\rho_0\to\rinfty$ at infinity is
\begin{equation}
  \rho_0(r,\theta)
  \;=\; \rinfty\,\exp\!\bigl(-\kappa\,\Phi_{\mathrm{Kerr}}(r,\theta)\bigr),
  \label{P7:eq:rho_Kerr_exact}
\end{equation}
where the exterior Newtonian potential associated with the Kerr
source is the multipole series
\begin{equation}
  \boxed{\Phi_{\mathrm{Kerr}}(r,\theta)
  \;=\; \frac{M}{r}
  \;-\; \frac{M a^2}{2r^3}\,P_2(\cos\theta)
  \;+\; \frac{3M a^4}{8r^5}\,P_4(\cos\theta)
  \;-\; \cdots
  \;=\; M\sum_{l=0}^{\infty}
        \frac{(-a^2)^l}{(2l-1)!!\,r^{2l+1}}\,P_{2l}(\cos\theta),}
  \label{P7:eq:Phi_Kerr}
\end{equation}
with $P_{2l}$ the Legendre polynomials and $(2l-1)!! = 1\cdot3\cdots(2l-1)$.
Each term satisfies $\nabla^2 r^{-(2l+1)}P_{2l} = 0$ individually,
so $\nabla^2\Phi_{\mathrm{Kerr}} = 0$ exactly term by term.
The series converges absolutely for $r > a$ (the exterior region
of any sub-extremal Kerr solution).

In the limit $a\to0$: $\Phi_{\mathrm{Kerr}}\to M/r$ and
$\rho_0\to\rinfty\,e^{-\kappa M/r}$, recovering the exact
Schwarzschild background.
\end{theorem}
\status{published}
\source{P7:thm:rotating}

\begin{remark}[Relation to the oblate spheroidal approximation]
\label{P7:rem:oblate}
The previously-used ansatz $1/\sqrt{r^2+a^2\cos^2\!\theta}$ is the generating
function of the Legendre polynomials:
\begin{equation}
  \frac{1}{\sqrt{r^2+a^2\cos^2\!\theta}}
  \;=\; \frac{1}{r}\sum_{l=0}^\infty
        \left(-\frac{a^2}{r^2}\right)^{\!l}
        \frac{(2l-1)!!}{(2l)!!}\,P_{2l}(\cos\theta),
  \label{P7:eq:oblate_expansion}
\end{equation}
which is \emph{not} term-by-term harmonic (each term has a different weight
from~\eqref{P7:eq:Phi_Kerr}), confirming that the oblate spheroidal
ansatz~\eqref{P7:eq:oblate_expansion} does not
satisfy $\nabla^2\rho_0=0$.  The exact background~\eqref{P7:eq:rho_Kerr_exact}
corrects this; numerically the two differ by less than $0.2\%$ for $a/r\lesssim0.5$.
\end{remark}
\status{published}
\source{P7:rem:oblate}

%════════════════════════════════════════════════════════════════════
% Kerr Hawking Temperature and the Weak Equivalence Principle — Paper VII
%════════════════════════════════════════════════════════════════════
The Hawking temperature from the acoustic surface gravity is
\begin{equation}
  T_H = \frac{\hbar\kappa_H}{2\pi k_B}
  \quad\text{(exact GR result at leading order)},
  \label{P7:eq:TH}
\end{equation}
with the soup correction from the sound-speed deviation $\delta c_s^2$:
\begin{equation}
  \boxed{T_H^{\mathrm{soup}} = T_H^{\mathrm{GR}}\!\left(1 + \frac{\beta^3}{4(1+\vp)}\right)
  \approx T_H^{\mathrm{GR}}\,(1+4.2\times10^{-5}).}
  \label{P7:eq:TH_correction}
\end{equation}
In the Schwarzschild limit $a\to0$, $T_H=\hbar c^3/(8\pi G_N M k_B)$
exactly; in the extremal limit $a\to r_s/2$, $T_H\to0$ exactly. The
correction $+4.2\times10^{-5}$ is currently unmeasurable but
constitutes a definite model-specific prediction.
\status{published}
\source{P7:sec:kerr}

After the one-loop induction of the Planck mass, all matter sectors
couple to the same induced metric $g_{\mu\nu}$; free-falling bodies
satisfy $Du^\mu/d\tau=0$ universally, regardless of composition, so the
WEP holds exactly at one-loop order via this geometric emergence. The
canonical field $\varphi_c$~\eqref{P7:eq:phi_c} satisfies a shift
symmetry
\begin{equation}
  \varphi_c \to \varphi_c + \delta\varphi_c
  \label{P7:eq:shift}
\end{equation}
of the kinetic sector (broken only by the potential). All couplings of
$\varphi_c$ to matter therefore appear through derivatives, which
couple universally to the stress tensor, so WEP violation from
derivative couplings is perturbatively suppressed:
\begin{equation}
  \eta_{\mathrm{pert}} \sim
  \left(\frac{m_\xi}{\LUV}\right)^2
  \sim \left(\frac{H_0}{\MPl}\right)^2
  \sim 10^{-122},
  \label{P7:eq:eta_pert}
\end{equation}
far below any experimental bound.
\status{published}
\source{P7:sec:wep}

The shift symmetry is broken non-perturbatively by instanton-like
field configurations that change the Hopf charge. For a test body
with $A$ baryons undergoing a coherent $\Delta H=1$ fluctuation, the
E\"otv\"os ratio is
\begin{equation}
  \eta_A \;=\; e^{-A\,S_1},
  \label{P7:eq:eta_A}
\end{equation}
where $S_1$ is the single-baryon instanton action derived from first
principles in Paper~V (Proposition~\ref{P5:prop:kink_avg}):
\begin{equation}
  \boxed{S_1
  \;=\; \frac{3}{4}\cdot 2^{-1/3}
        \cdot\left[\Bigl(\frac{5}{4}\Bigr)^{3/4}-1\right]
  \;\approx\; 0.1084.}
  \label{P7:eq:S1_FP}
\end{equation}
For a single baryon $\eta_1=e^{-0.108}\approx0.90$ (WEP maximally
violated at the nucleon scale); for $A\gtrsim277$ baryons $\eta<10^{-13}$
(MICROSCOPE bound satisfied); for macroscopic bodies ($A\sim10^{26}$):
\begin{equation}
  \eta \;\sim\; e^{-0.1084\times10^{26}} \;\approx\; 10^{-10^{25}},
  \label{P7:eq:eta_macro}
\end{equation}
the WEP holds to effectively infinite precision.
\status{published}
\source{P7:sec:wep}

